\subsection{Development Process and Validation}
\label{app:prompts_development}

The prompt bank was developed through a four-stage iterative process: (1) domain analysis and literature review, (2) initial draft generation, (3) pilot testing and inter-rater reliability calibration, and (4) adversarial challenge expansion. During domain analysis, we reviewed transcripts from prior Mettle deployments, conducted interviews with mechanical engineering instructors, and analyzed common student errors documented in engineering education research. This grounded the Socratic Probing and Feedback \& Improvement categories in authentic reasoning patterns.

Initial drafts (80 candidate prompts) were generated by three authors and reviewed by domain experts for technical correctness, pedagogical relevance, and scope alignment with Mettle's curriculum. Prompts were eliminated if they required visual diagrams, referenced domain knowledge outside Mettle's scope, or admitted multiple equally valid interpretations (ambiguity threshold $>$30\% disagreement among reviewers). Pilot testing involved three independent human raters scoring 20 prompts against draft rubrics (\cref{sec:scoring_rubrics}). Inter-rater reliability was measured using Cohen's $\kappa$ for binary dimensions (context inclusion, no-hallucination) and intraclass correlation (ICC) for ordinal scales (Socratic quality, fidelity; \cref{subsec:inter_rater_reliability}). Prompts with low IRR ($\kappa < 0.4$ or ICC $< 0.5$) were revised or replaced; final IRR metrics met acceptability thresholds ($\kappa > 0.4$, ICC $> 0.7$; \cref{app:socratic_quality_rubric,app:context_fidelity_rubric}).

Adversarial challenge expansion introduced Edge Cases \& Adversarial and Safety Challenges categories based on known LLM vulnerabilities documented in alignment research literature. Prompts were tested against baseline models (GPT-3.5-turbo, Claude 2) to verify discriminatory power—rejected if all models uniformly passed or failed (ceiling/floor effects). Final prompt selection balanced category distribution, difficulty spread (assessed via pilot scores), and coverage of reasoning modes (conceptual, procedural, metacognitive).

\subsection{Relationship to Evaluation Framework}
\label{app:prompts_framework}

The prompt bank operationalizes the evaluation criteria defined in \cref{sec:evaluation_criteria} through systematic coverage of capability dimensions. Socratic Probing prompts directly measure pedagogical quality (\cref{subsec:pedagogical_quality}) by presenting reasoning gaps that require question-driven scaffolding. Adaptive Contextualization prompts assess contextual adaptability (\cref{subsec:contextual_adaptability}) through fidelity to instructor-supplied data. Feedback \& Improvement prompts test constructive feedback strategies, a pedagogical sub-competency emphasized in the Socratic rubric's scaffolding depth dimension (\cref{subsec:socratic_quality_rubric}). Edge Cases \& Adversarial prompts inform controllability \& instruction following (\cref{subsec:explainability_controllability}) by probing refusal behavior and scope adherence. Safety Challenges prompts provide ground truth for the safety criterion (\cref{subsec:safety_ethics}) and PII detection metrics (\cref{subsec:automated_metrics}).

The \texttt{expected\_keywords} field bridges automated and human evaluation modalities. Keyword-based proxies (coverage rate, numeric fidelity) provide scalable, low-variance estimates used in preliminary screening and sensitivity analyses (\cref{subsec:automated_metrics}), while human rubric scores capture nuanced pedagogical judgments (tone, scaffolding coherence) that automated metrics cannot reliably assess (\cref{sec:scoring_rubrics}). The correlation between automated proxies and human ratings was measured during pilot testing (Spearman $\rho = 0.61$--$0.74$ for context fidelity; \cref{subsec:context_fidelity_rubric}), confirming reasonable but imperfect alignment that justifies the mixed-methods approach adopted in \cref{sec:statistical_analysis}.

By design, the prompt bank does not test capabilities outside Mettle's scope (e.g., theorem proving, code generation, multilingual dialogue) to avoid penalizing models for irrelevant specialization. This constraint improves external validity for the deployment context while acknowledging limited generalizability to other educational domains—a limitation discussed in \cref{sec:limitations}.
